\documentclass[11pt,a4paper]{article}
\usepackage[margin=24mm]{geometry}
\usepackage[T1]{fontenc}
\usepackage{lmodern}
\usepackage{amsmath,amssymb,booktabs,longtable,array}
\usepackage{enumitem,microtype}
\usepackage[hidelinks]{hyperref}
\urlstyle{same}
\setlist{nosep,leftmargin=*}
\setlength{\emergencystretch}{2em}
\newcommand{\sex}{s_{\mathrm{ex}}}
\newcommand{\Hv}{H_{\mathrm V}}
\newcommand{\study}{\path{ComponentEntropyStudy_20261030/}\allowbreak}
\title{Voronoi and Graph Entropy in Multicomponent Charged Liquids\\
\large Work plan for an independent structural assessment}
\author{EntropicFluids collaboration}
\date{31 October 2026}
\begin{document}
\maketitle
\begin{abstract}
This plan proposes a staged analysis of existing molecular-dynamics trajectories:
first local Voronoi topology and chemical neighborhoods, then, if sampling permits,
cluster-entropy extrapolation. Thermodynamic integration (TI) remains the reference
for excess entropy. The objectives are to determine whether structural diversity
tracks TI entropy, whether environments around unchanged divalent ions explain
their mobility, and whether adding physically distinct components has effects
beyond adding species labels. No additional MD is required for the initial pilot.
The protocol below is proposed work, not a report of completed Voronoi calculations.
\end{abstract}

\section{Scientific questions and prior work}
Rosenfeld scaling relates reduced transport to excess entropy~\cite{Rosenfeld1977}.
Charged-mixture applications precede the present study~\cite{Rosenfeld1995}.
Generalized tracer scaling has been tested in model mixtures, with limitations
when component mobilities decouple~\cite{Krekelberg2009}. More recent mixture
modeling uses entropy together with limiting transport properties and mixing
rules; one bulk entropy is not generally a unique predictor of each diffusion
coefficient~\cite{Schmitt2025}.

For electrolytes, Wang et al. connect multiple salts, altered solvation, and
enhanced lithium mobility~\cite{Wang2023}. He et al. explicitly report a
Rosenfeld relation for lithium self-diffusion in multi-anion
electrolytes~\cite{He2026}. These motivate comparison, but do not establish a
universal causal effect of component count. Only the abstract and bibliographic
record of Ref.~\cite{He2026} were accessible during preparation of this plan;
its entropy estimator and controls must be checked in the full paper before
making a detailed methodological comparison.

Two graph approaches should be distinguished:
\begin{enumerate}
\item \textbf{Local cell topology.} Yoon et al.~\cite{Yoon2019} classify Voronoi
cells by their complete graph topology, using Weinberg vectors, and compare
topological Shannon entropy with TI entropy and transport in repulsive
soft-sphere fluids. Their entropy conversion is interaction dependent; rare
types and correlations between adjacent cells matter. This is the closest
precedent for the initial pilot. The underlying topology framework and an
implementation are described in Refs.~\cite{Lazar2015,LazarToolkit}.
\item \textbf{Cluster entropy.} Zhou and Milner~\cite{Zhou2016} classify
Voronoi-derived clusters by graph isomorphism and estimate structural entropy
per particle from the growth of subsystem entropy with particle number.
They benchmark hard-particle systems against previous configurational-entropy
estimates from TI. This motivates the more expensive second stage, but does
not prove equivalence to excess entropy in our attractive, charged mixtures.
\end{enumerate}
Voronoi entropy has also been used for hydrogen-bonded liquids and binary
mixtures~\cite{Abdulkareem2024}; that application illustrates structural
characterization rather than a general absolute-entropy prescription.

The proposed contribution is a benchmark across controlled charged mixtures,
with an exact label control and independent TI and diffusion data. The literature
search establishes precedent, not priority for every element of this design.

\section{Entropy definitions and comparison rules}
The reference is classical excess entropy per particle relative to the ideal
mixture at the same temperature, density, and composition:
\begin{equation}
\sex = \frac{S-S_{\mathrm{id,mix}}}{N},\qquad
\frac{\sex}{k_B}=\frac{U_{\mathrm{ex}}}{Nk_BT}
-\frac{A_{\mathrm{ex}}}{Nk_BT}.
\end{equation}
Ideal mixing, $-k_B\sum_i x_i\ln x_i$, is excluded. Do not substitute the
excess entropy of mixing relative to pure liquids for this quantity.

For a chosen local topological class $\alpha$, define the dimensionless
information entropy
\begin{equation}
\Hv=-\sum_\alpha p_\alpha\ln p_\alpha.
\end{equation}
This describes the diversity of a coarse structural representation. It is not
automatically $\sex/k_B$, and summing single-cell entropies omits correlations.
Likewise, glass configurational entropy, which counts basins after separating
vibrational motion, is a different target from our full liquid excess entropy.

Use a simulated ideal-gas point ensemble with the same periodic geometry and
the identical estimator to define an optional topological reference
$\Delta H_{\mathrm V}=\Hv-H_{\mathrm{V,id}}$. Do not assume that subtracting
this reference makes it a thermodynamic excess entropy. If chemical colors
are used, match their composition in the reference. Any fitted mapping to TI
must be explicitly called a calibration and tested on excluded compositions.

The completed classical 2PT assessment follows
Refs.~\cite{Sun2017,Lai2012}. Its internal checks passed, but its absolute
bias varies with composition and it overestimates one controlled TI increment.
Keep it as a secondary comparator. In particular, diffusion enters its spectral
partition, whereas the proposed static graph descriptors do not use velocities.

\section{Systems, existing evidence, and priorities}
All states have $N=4000$, $T^*=2$, $\rho^*=0.7$, $x_{+2}=x_{-2}=0.03$,
unit masses and nominal diameters, and three replicas. Ions and primary neutral
solvent have $\varepsilon=1$. Unlike LJ amplitudes follow geometric mixing.
The shifted LJ cutoff is fixed at $2.5\sigma$ throughout TI; electrostatics use
PPPM at tolerance $10^{-5}$. The parameter $I=\tfrac12\sum_i x_iq_i^2$ below
is dimensionless; it is not a molar ionic strength.

\begin{table}[ht]
\centering\small
\caption{Composition map. Secondary fractions and amplitudes correspond in order.
The monovalent fraction is for each sign; all rows also contain 3\% of each divalent sign.}
\begin{tabular}{@{}lccccc@{}}
\toprule
Code & Primary $x$ & Secondary $x$ & Secondary $\varepsilon$ & $x_{\pm1}$ & $I$\\
\midrule
B3 & .94 & --- & --- & 0 & .12\\
B5 & .84 & --- & --- & .05 & .17\\
U4 & .64 & .30 & .25 & 0 & .12\\
L5 & .64 & .15+.15 & .25, .25 & 0 & .12\\
M5 & .64 & .15+.15 & .0625, .5625 & 0 & .12\\
H5 & .64 & .15+.15 & .01, .81 & 0 & .12\\
U6 & .54 & .30 & .25 & .05 & .17\\
H7 & .54 & .15+.15 & .01, .81 & .05 & .17\\
\bottomrule
\end{tabular}
\end{table}

\begin{enumerate}
\item \textbf{Mandatory null: U4/L5.} Identical interactions, with one secondary
solvent split into two labels. Formal ideal mixing changes by
$0.30\ln2=0.20794$ per particle in units of $k_B$; diffusion is identical in
the completed trajectories. TI differences are consistent with sampling noise.
Verify positional identity by ID rather than assuming it from equal diffusion.
Also create a synthetic label split of the exact same snapshots to guarantee
an exact implementation test.
\item \textbf{Primary contrasts: H5--U4 and H7--U6.} Both TI increments are
approximately $+0.1033k_B$. Primary-window probe diffusion increases by
4.81\% and 5.20\%, respectively. The former loses nominal significance in
the late MSD window; the latter is unresolved across replicas. These small
effects must not be described as established merely because a structural
descriptor changes significantly.
\item \textbf{Moderate perturbation: M5--U4.} TI increases by $0.0304k_B$;
the approximately 1.5\% probe response is unresolved. Include this case to
avoid selecting only large structural changes.
\item \textbf{Broad references: B3/U4/H5 and B5/U6/H7.} Baseline-to-design
probe increases are about 20\% and 18\%. These comparisons change mean
interaction strength and supply range, rather than an isolated component-count
test. B3/B5 also changes ionic strength and is a secondary contrast.
\end{enumerate}

U4/L5/M5/H5 have identical $\sum_i x_i\sqrt{\varepsilon_i}=0.85$ and mean
bare pair amplitude $\sum_{ij}x_ix_j\varepsilon_{ij}=0.7225$; U6/H7 also
share these values. Actual energy, effective excluded volume, solvation, and
friction need not match. L5 is a duplicate control, not another independent
state in regressions. The pilot should cover all eight states if processing
is inexpensive; otherwise begin with U4/L5/H5/U6/H7 and B3, then add M5/B5
before drawing study-wide conclusions.

\section{Data handover and reproducibility}
Paths are relative to the repository root; raw trajectories remain local and
are not included in the Git repository. Supply them separately to the analyst.
\begin{itemize}
\item \study\texttt{selected.json}: exact species definitions and case names.
\item \study\texttt{md/<case>/replica\_<r>/transport/positions.nc}:
all-particle, unwrapped positions every $0.5\tau$ over $1000\tau$ production.
Read stored IDs, types, frame times, and box information explicitly; do not
infer atom ordering from row number.
\item Corresponding \texttt{transport/data.mixture}, \texttt{in.run}, and
replica \texttt{manifest.json}: charges/type mapping, box, inputs, and provenance.
\item \texttt{results/summary.json}, \texttt{results/diagnostics.json}, and
\texttt{results/conclusions.md} under \study: TI and transport comparisons,
intervals, estimator sensitivities, and interpretation.
\item \texttt{md/<case>/replica\_<r>/replica\_result.json} and
\texttt{transport/msd\_multi\_origin.csv}: replica-level reference values.
\item \study\texttt{two\_pt\_rehabilitation/}: corrected 2PT methods,
replica results, \texttt{summary.json}, and \texttt{report.md}.
\end{itemize}
Case directories are B3: \path{baseline_3C}; B5: \path{mono_f010};
U4: \path{weak_f030_e025}; L5: \path{dispersion_split_d000};
M5: \path{dispersion_split_d025}; H5: \path{dispersion_split_d040};
U6: \path{mono_weak_uniform6C}; H7: \path{mono_weak_split7C}.

Record source commit, checksums, tool versions, all tessellation/graph options,
selected frame IDs, random seeds, and estimator settings. Retain canonical
graph identifiers and type counts, not just entropy estimates. Write new
outputs to a separate directory such as \texttt{VoronoiEntropyAssessment/}.
Use the colleague's validated tools; the scientific definitions below should
be agreed before comparing descriptor values with diffusion.

\section{Stage A: local structural pilot}
\subsection{Geometry and graph validation}
Wrap coordinates into each frame's periodic box before tessellation; preserve
the unwrapped record for displacement analysis. Tessellate \emph{all} particles,
then condition statistics on the central species. Tessellating only ions would
measure a different neighborhood and is not the primary protocol. Equal nominal
diameters motivate ordinary Voronoi cells initially. Weighted cells, if used
later, require an explicit radius convention and sensitivity analysis.

Validate periodic cells by checking that volumes sum to the box volume,
adjacency is reciprocal, and results are invariant under translation, atom
ordering, and type-code renaming. Distinguish the edge graph of a single
polyhedron from a particle-adjacency graph. Canonicalization must resolve graph
isomorphism exactly; a hash or face-count vector alone is not a complete
topological identifier. Check any treatment of tiny faces or near-degenerate
vertices on a fixed subset before selecting a primary convention.

Use instantaneous configurations as primary. Quenching is an optional,
separately identified inherent-structure analysis; it changes the target and
must not silently replace liquid configurations.

\subsection{Descriptors and controls}
Produce the following for the full liquid and separately for $+2$, $-2$, and
each mechanically distinct neutral species:
\begin{enumerate}
\item Full uncolored cell-topology frequencies and $\Hv$; retain face-count
and coordination distributions as simpler diagnostic summaries.
\item Cell volume, surface area, and shape distributions. These capture metric
information discarded by topology; report their moments without treating a
bin-dependent differential entropy as an absolute thermodynamic entropy.
\item Neighbor composition and charge distributions around the central ion.
If supported, calculate colored neighborhood graph classes. Specify precisely
whether colors belong to faces (neighbor particles) or nodes of a particle
graph, and canonicalize colors consistently with graph automorphisms.
\item Geometry-only, charge-only, and full mechanical-species versions, with
mechanically identical L5 labels merged in the physical-species version.
Also report label-resolved L5 as a deliberate bookkeeping diagnostic.
\end{enumerate}

At fixed coordinates, permute species labels while preserving counts and
holding the selected central ion identity fixed for ion-centered statistics.
Use multiple seeded permutations, preferably at least 100 for the pilot.
This is a chemical-association null, not an equilibrium realization of the
charged Hamiltonian. Compare actual and randomized conditional neighborhood
distributions to expose ordering beyond composition. Keep purely combinatorial
color diversity separate from geometry. Entropy differences against permutation
controls are structural statistics, not automatically thermodynamic contributions.

\subsection{Sampling and uncertainty}
Start with 100 frames evenly spaced over each dense production trajectory;
repeat with 200 and 400 frames if available and needed. Increase separation
or use blocks after measuring correlation of representative topology and
coordination observables. All 4000 particles can contribute to bulk statistics,
but only 120 ions per sign occur per frame and neighboring cells are correlated.

Report raw entropy, bias-corrected estimates supported by the tools, singleton
fractions, discovery curves, and rare-type coverage. An unseen-type estimator
does not create independent samples. Repeat with matched sample sizes and
thinning. If entropy continues to grow with sample count, report an unresolved
estimate rather than selecting the apparent plateau. Include independent
ideal-gas point ensembles with the same pipeline; their topology can require
much more sampling than the dense liquid.

Keep replicas separate throughout estimation. Use paired differences across
the three replica indices for composition contrasts; report nominal 95\%
Student-$t$ intervals with two degrees of freedom. Block resampling within
replicas diagnoses temporal uncertainty but does not turn cells, overlapping
clusters, frames, or permutations into independent MD replicas. Check first
and second trajectory halves for stationarity. Monitor cation/anion symmetry;
the existing baseline transport asymmetry is a sampling flag.

\section{Stage B: thermodynamic and transport assessment}
Compare descriptor \emph{increments} first, particularly H5--U4 and H7--U6.
Then compare absolute-state ordering against TI, HNC, and corrected 2PT.
Calculate the independent pair contribution from species RDFs where feasible:
\begin{equation}
\frac{s_2}{k_B}=-2\pi\rho\sum_{ij}x_ix_j
\int_0^\infty [g_{ij}(r)\ln g_{ij}(r)-g_{ij}(r)+1]r^2\,dr.
\end{equation}
The sum is over ordered species pairs. Document finite-box cutoff, binning,
normalization, and tail sensitivity; merging identical labels must preserve
the result. This gives a simple structural baseline before crediting graph
analysis with additional predictive information.

For transport use the existing common-system-COM MSD estimates:
\begin{equation}
D_{\mathrm{net}}=\sum_i x_iD_i,\quad
D_{\mathrm{probe}}=\tfrac12(D_{+2}+D_{-2}),\quad
D_i^{\mathrm R}=D_i\rho^{1/3}\sqrt{m_i/(k_BT)}.
\end{equation}
Use total number density in the reduction. At common $T$, $\rho$, and mass,
diffusion ratios are unchanged by it. Neither $D_{\mathrm{net}}$ nor
$D_{\mathrm{probe}}$ is conductivity, mutual diffusion, or diffusion of a
bound neutral ion pair. Retain the primary 20--100$\tau$ MSD fit and the
10--50 and 50--200$\tau$ sensitivity windows.

Compare $\ln D_i^{\mathrm R}$ with bulk topology and ion-centered descriptors.
Keep fixed-$I$ families separate; exclude L5 from fit counts. With only seven
physical states, a high training $R^2$ is weak evidence. Limit models to one
descriptor plus an intercept initially; show residuals, paired contrasts,
leave-one-state-out errors, and transfer between the two salt backgrounds.
Do not select descriptors or graph parameters by maximizing agreement with
diffusion and then present that agreement as an independent test.

As an optional mechanistic extension, relate an ion's initial environment to
its subsequent squared displacement at fixed lags (e.g., 5, 20, and $50\tau$),
with the same whole-system COM correction as the reference analysis. Call
this conditional mobility, not a local self-diffusion coefficient. Analyze
within each state so that between-state differences cannot create a spurious
local association. Neighbor persistence may help distinguish geometric packing
from chemical residence effects; the $0.5\tau$ output spacing limits short-time
transition and residence-time resolution.

\section{Stage C: conditional cluster-entropy extrapolation}
Proceed only after Stage A passes the label and convergence checks and the
analyst demonstrates adequate graph-count coverage. First reproduce a small
hard-particle benchmark using the cluster construction and entropy convention
of Zhou and Milner~\cite{Zhou2016}. Their full method must be consulted for
cluster selection, boundary treatment, and corrections; the schematic formula
below is not a replacement implementation.

For a fixed, documented rule for selecting connected $n$-particle neighborhoods,
estimate canonical graph-class probabilities and
\begin{equation}
H_G(n)=-\sum_g p_n(g)\ln p_n(g),\qquad
h_G(n)=H_G(n+1)-H_G(n).
\end{equation}
Begin with an adaptive sequence such as $n=2,3,4,5,6,8,10$, extending only
while class coverage remains credible. Rooted and unrooted graphs, graph-radius
and nearest-$n$ neighborhoods, and induced versus boundary-completed graphs
define different distributions: select one primary convention and document
all alternatives. Preserve periodic adjacency before extracting local clusters.
Distinguish clusters sampled around probes from bulk clusters; their entropy
slopes have different conditioning.

Evaluate uncolored graphs first, followed by charge/species colors with merged
identical labels. Apply the same ideal-gas and fixed-composition permutation
references. Check sample-size growth for every $n$: saturation of observed
entropy at roughly the logarithm of the number of sampled clusters is a
coverage failure, not physical convergence.

A defensible extensive regime would have stable $h_G(n)$ or stable fitted
slopes under changes of fit range, trajectory half, sample size, and replica.
Test boundary corrections rather than presuming exact linearity for small
clusters. Require at least four consecutive, adequately sampled integer sizes
for a provisional slope fit. If a converged regime cannot be reached before
rare graphs dominate, stop and report finite-cluster information statistics.
Do not force an infinite-size extrapolation.

Even a stable slope is entropy of the chosen graph representation. In an
energetic liquid, configurations with the same graph can have different
distances, energies, and vibrational phase-space volumes. The missing
within-class information must be accounted for, or the result retained as
structural entropy. A TI calibration is an empirical mapping and requires
held-out tests. It does not establish an exact statistical-mechanical identity.

\section{Decisions, deliverables, and follow-up}
\begin{longtable}{@{}p{0.25\textwidth}p{0.69\textwidth}@{}}
\toprule
Checkpoint & Required evidence and action\\
\midrule
\endhead
Implementation & Periodicity and invariance checks pass; a synthetic label split
leaves geometry and merged-color results identical to numerical precision.
Audit actual U4/L5 snapshots before interpreting any discrepancy.\\[3pt]
Sampling & Discovery curves, thinning, halves, and replica analyses are stable
at the scale of the target contrasts. Otherwise extend postprocessing samples
or report unresolved descriptors; add MD only if available sampling is inadequate.\\[3pt]
Thermodynamic proxy & Report raw structural results independently. If a
calibrated $\sex$ predictor is proposed, use provisional practical targets
of $0.02k_B$ error in controlled increments and $0.05k_B$ in absolute values,
with held-out validation. These are project targets, not literature guarantees;
wide intervals do not establish equivalence.\\[3pt]
Mechanistic value & Establish whether charge/species-aware probe descriptors
explain behavior beyond bulk topology, TI, or $s_2$. An unresolved 5\% transport
contrast remains unresolved even if its structural contrast is clear.\\[3pt]
Cluster extension & Attempt extrapolation only with validated construction,
adequate class coverage, and a stable entropy-growth regime. No stable regime
means no reported thermodynamic entropy from extrapolation.\\
\bottomrule
\end{longtable}

The handover should contain a machine-readable per-state/per-replica table;
paired contrast tables with intervals; graph and chemical-null definitions;
sample-coverage and estimator-sensitivity plots; TI--topology and
diffusion--descriptor comparisons; and, if attempted, $H_G(n)$, increments,
fit-range sensitivity, and explicit extrapolation failures. Include scripts,
tool versions, frame manifests, and enough graph frequencies to reproduce
the entropy estimates. A short written assessment should distinguish
thermodynamic prediction, structural association, and mechanistic interpretation.

The initial effort is postprocessing of existing trajectories. Time the first
few frames and graph sizes before estimating total CPU/memory cost. Local
cell classification and chemical-null analysis come first; cluster enumeration
is the potentially expensive part. No new GPU submission is part of this plan.
If structure is well resolved but transport is not, prioritize additional
replicas, longer dense trajectories, and a second box size at fixed physical
state. Higher concentrations, charge asymmetry, and different nominal sizes
are later transfer tests; size changes also require a documented ordinary
versus weighted Voronoi comparison.

\begin{thebibliography}{99}
\small
\bibitem{Rosenfeld1977}
Y. Rosenfeld, ``Relation between the transport coefficients and the internal
entropy of simple systems,'' \textit{Phys. Rev. A} \textbf{15}, 2545--2549 (1977).
\href{https://doi.org/10.1103/PhysRevA.15.2545}{doi:10.1103/PhysRevA.15.2545}.

\bibitem{Rosenfeld1995}
Y. Rosenfeld, E. Nardi, and Z. Zinamon, ``Corresponding States Hard-Sphere Model
for the Diffusion Coefficients of Binary Dense-Plasma Mixtures,''
\textit{Phys. Rev. Lett.} \textbf{75}, 2490 (1995).
\href{https://doi.org/10.1103/PhysRevLett.75.2490}{doi:10.1103/PhysRevLett.75.2490}.

\bibitem{Krekelberg2009}
W. P. Krekelberg, M. J. Pond, G. Goel, V. K. Shen, J. R. Errington, and
T. M. Truskett, ``Generalized Rosenfeld scalings for tracer diffusivities in
not-so-simple fluids: Mixtures and soft particles,'' \textit{Phys. Rev. E}
\textbf{80}, 061205 (2009).
\href{https://doi.org/10.1103/PhysRevE.80.061205}{doi:10.1103/PhysRevE.80.061205};
\href{https://arxiv.org/abs/0910.0280}{accessible preprint}.

\bibitem{Schmitt2025}
S. Schmitt, H. Hasse, and S. Stephan, ``Entropy scaling for diffusion
coefficients in fluid mixtures,'' \textit{Nature Communications}
\textbf{16}, 2611 (2025).
\href{https://doi.org/10.1038/s41467-025-57780-z}{doi:10.1038/s41467-025-57780-z}.

\bibitem{Wang2023}
Q. Wang, C. Zhao, J. Wang, Z. Yao, S. Wang, S. G. H. Kumar, S. Ganapathy,
S. Eustace, X. Bai, B. Li, and M. Wagemaker,
``High entropy liquid electrolytes for lithium batteries,''
\textit{Nature Communications} \textbf{14}, 440 (2023).
\href{https://doi.org/10.1038/s41467-023-36075-1}{doi:10.1038/s41467-023-36075-1}.

\bibitem{He2026}
L. He, Y. Lin, J. Li, H. Su, H. Liu, C. Yang, and C. Lian,
``Bridging thermodynamics and ion transport in lithium batteries via molecular
design of high-entropy electrolytes,'' \textit{Molecular Systems Design \&
Engineering} \textbf{11}, 681--690 (2026).
\href{https://doi.org/10.1039/D6ME00060F}{doi:10.1039/D6ME00060F}.
First published 19 May 2026; abstract-level comparison in this plan.

\bibitem{Yoon2019}
T. J. Yoon, M. Y. Ha, E. A. Lazar, W. B. Lee, and Y.-W. Lee,
``Topological extension of the isomorph theory based on the Shannon entropy,''
\textit{Phys. Rev. E} \textbf{100}, 012118 (2019).
\href{https://doi.org/10.1103/PhysRevE.100.012118}{doi:10.1103/PhysRevE.100.012118};
\href{https://arxiv.org/abs/1901.02772}{accessible preprint}.

\bibitem{Lazar2015}
E. A. Lazar, J. Han, and D. J. Srolovitz, ``A topological framework for local
structure analysis in condensed matter,'' \textit{Proceedings of the National
Academy of Sciences USA} \textbf{112}, E5769--E5776 (2015).
\href{https://doi.org/10.1073/pnas.1505788112}{doi:10.1073/pnas.1505788112};
\href{https://arxiv.org/abs/1508.05937}{accessible preprint}.

\bibitem{LazarToolkit}
E. A. Lazar, ``VoroTop: Voronoi Cell Topology Visualization and Analysis
Toolkit,'' \textit{Modelling and Simulation in Materials Science and
Engineering} \textbf{26}, 015011 (2018; published online December 2017).
\href{https://doi.org/10.1088/1361-651X/aa9a01}{doi:10.1088/1361-651X/aa9a01};
\href{https://arxiv.org/abs/1804.04221}{accessible preprint};
\href{https://www.vorotop.org/}{software and documentation}.

\bibitem{Zhou2016}
Y. Zhou and S. T. Milner, ``Structural entropy of glassy systems from graph
isomorphism,'' \textit{Soft Matter} \textbf{12}, 7281--7288 (2016).
\href{https://doi.org/10.1039/C6SM01355D}{doi:10.1039/C6SM01355D}.
Primary source for the cluster-entropy extension.

\bibitem{Abdulkareem2024}
Abdulkareem U and V. Madhurima, ``Understanding complexity of hydrogen-bonded
liquids through Voronoi entropy,'' \textit{Chemical Physics Impact}
\textbf{8}, 100403 (2024; available online 30 November 2023).
\href{https://doi.org/10.1016/j.chphi.2023.100403}{doi:10.1016/j.chphi.2023.100403}.

\bibitem{Sun2017}
T. Sun, J. Xian, H. Zhang, Z. Zhang, and Y. Zhang,
``Two-phase thermodynamic model for computing entropies of liquids reanalyzed,''
\textit{J. Chem. Phys.} \textbf{147}, 194505 (2017).
\href{https://doi.org/10.1063/1.5001798}{doi:10.1063/1.5001798};
\href{https://arxiv.org/abs/1711.02041}{accessible preprint}.

\bibitem{Lai2012}
P.-K. Lai, C.-M. Hsieh, and S.-T. Lin,
``Rapid determination of entropy and free energy of mixtures from molecular
dynamics simulations with the two-phase thermodynamic model,''
\textit{Phys. Chem. Chem. Phys.} \textbf{14}, 15206--15213 (2012).
\href{https://doi.org/10.1039/C2CP42011B}{doi:10.1039/C2CP42011B}.
\end{thebibliography}
\end{document}
